% Active Case F certificate: fixed-start points and explicit comparisons.
% This source is shared verbatim by the canonical and immediate-proof editions.
The frontier witnesses remain $Q_-,Q_0,Q_+$.  The comparison points $A,B,C$
are the fixed-start Newton points of
\zcref{lem:technical-newton-reduction}; they are not additional forced
frontier witnesses.  The actual disk has radius $\eta=h(1-c_*)$.
Only the tangent-residual calculation uses a smaller scalar radius.
The construction below replaces the former sparse Bernstein transcript,
not the geometric forcing or the two straight-line contact estimates.
It remains an exact computer-assisted argument.
The five new comparison checks use $383$ Chebyshev coefficients, the curvature
comparison uses $113$, and the two increment comparisons use $208$ nonconstant
coefficients; the boundary Taylor and identity checks are additional.
The historical junction-start transcript and its two exact programs remain
unchanged regression material, not inputs for these new point formulas.


\Fcertificateheading{A single rational radius envelope}
Reflect so that $a\ge b$ and set
\[
 \begin{gathered}
 m=1-a,\quad U=a+b-1,\quad a=1-m,\quad b=m+U,\\
 \omega=m(1-m),\qquad b_0=\omega/2.
 \end{gathered}
\]
Here $U$ is a scalar, not an open triangle $U_i$.
The quartic capacity root $C_L(m)$ was defined in
\zcref{eq:core-cstar}.  We first obtain a rational lower bound for
$1-C_L(m)$, then extend it to both capacity branches.

\begin{lemma}[A rational upper-start Newton deficit]
\label{lem:f-rational-newton-deficit}
Let $0<m<1/2$, $h=\sqrt3/2$, and let $C_L(m)$ be the unique root in
$[h,1]$ of $F_m(c)=c^4-c^2+mc-m^2$.
Set $\omega=m(1-m)$ and $c_0=1-\omega/2$. Then
\[
 r(m):=1-c_0+\frac{F_m(c_0)}{F_m'(c_0)}
 =\frac{\omega(4-\omega)(3\omega^2-4\omega+4)}
 {8(4+2m-10\omega+6\omega^2-\omega^3)}
\]
satisfies
\[
0<\omega/2<r(m)\le1-C_L(m)<m.
\]
\end{lemma}
\begin{proof}
One has $c_0\ge7/8>h$ and
\[
 F_m(c_0)=\frac{m^2(1-m)}{16}P(m),\qquad
 P(m)=12-28m+17m^2-11m^3+3m^4-m^5.
\]
On $[0,1/2]$, $P'\le-28+17+3/2=-19/2$ and
$P(1/2)=33/32>0$. Thus $F_m(c_0)>0$.
The polynomial $F_m$ is increasing and strictly convex on $[h,1]$,
so $c_0>C_L(m)$ and its tangent zero
$c_1=c_0-F_m(c_0)/F_m'(c_0)$ satisfies
$C_L(m)\le c_1<c_0$. This proves the first bounds.
If $1-m<h$, then $C_L(m)>1-m$ immediately. Otherwise it follows from
\[
 F_m(1-m)=-m((1-m)^3+m^2)<0.
\]
The displayed rational formula follows by substitution; its denominator
factor $4+2m-10\omega+6\omega^2-\omega^3=2F_m'(c_0)$ is positive.
\end{proof}

\begin{lemma}[Uniform affine radius envelope]
\label{lem:f-uniform-affine-radius-envelope}
\label{lem:paper-branchwise-cbar}
In the strict F domain, put $U=a+b-1$, $m=\min(1-a,1-b)$, and let
$c_*$ be the exact common-pair capacity. For any
$0<r\le1-C_L(m)$, define
\[
 E_r(m,U)=r+(m-r)\left(1-\frac Ur\right)_+
 =\max\left\{r,m-\frac{m-r}{r}U\right\}.
\]
Then
\[
 1-c_*\ge E_r(m,U).
\]
In particular, the rational function in
Lemma~\ref{lem:f-rational-newton-deficit} gives one explicit radius recipe
without testing the original capacity-branch selector.
\end{lemma}
\begin{proof}
Reflect to $a\ge b$ and set $s=1-U$, $C=C_L(m)$. The exact capacity
formula is $c_*=C$ on the quartic branch and
\[
 g_m(s)=\frac{2(s-m)}{1+\sqrt{4s^2-3}}
\]
on the triangular branch. On the latter, its selector
$\chi=F_m(s)>0$ gives $s>h$ because
$\chi\le s^4-3s^2/4$, and hence $s>C$.
Writing $E=\sqrt{4s^2-3}$,
\[
 g_m'(s)=\frac{2(E-3+4ms)}{E(1+E)^2}\le0
 \qquad(C<s\le1).
\]
The root equation gives
$m/C=(1\pm\sqrt{4C^2-3})/2$, and substitution shows $g_m(C)\le C$.
Thus $c_*\le C$ on both branches, so $1-c_*\ge r$.

It remains to consider $0<U<r$. Here $s>1-r\ge C$, so the triangular
formula applies. On the analytic interval $s\in[1-r,1]$,
\[
 g_m''(s)=\frac{8\{s(3+9E-2E^2)-m(3+9E+2E^3)\}}
 {E^3(1+E)^3}>0.
\]
Indeed $m/s\le1/\sqrt3<2/3$ and
\[
 (3+9E-2E^2)-\frac23(3+9E+2E^3)
 =1+3E-2E^2-\frac43E^3\ge\frac23
 \quad(0\le E\le1),
\]
since the last cubic is concave and lies above its endpoint chord.
Endpoint cases follow by continuity. Therefore $U\mapsto g_m(1-U)$
is convex. Its endpoint bounds are
$g_m(1)=1-m$ and $g_m(1-r)\le C\le1-r$.
The chord inequality gives
\[
 c_*\le(1-U/r)(1-m)+(U/r)(1-r),
\]
which is the claimed affine bound. Combining it with $1-c_*\ge r$
proves the result.
\end{proof}


For this choice of $r$ put
\begin{equation}
 \kappa=\frac{m-r}{r},\qquad
 e_0=\max\{r,m-\kappa U\},\qquad e_*=1-c_*.
 \label{eq:f-uniform-radius}
\end{equation}
Thus $0<b_0\le r\le e_0\le e_*$, $r<m$, and $r\le1-h<1/6$.
The slope satisfies
\begin{equation}
 1+\frac m2+2m^2\le\kappa\le3,
 \qquad \ell(m):=\frac{1+5m}{2}\le\kappa.
 \label{eq:f-slope-bounds}
\end{equation}
For the upper bound use $r\ge m(1-m)/2$.  For the lower bound write
$r=N_r/L_r$, using the positive numerator and denominator in
\zcref{lem:f-rational-newton-deficit}.  Exact expansion gives
\[
 mL_r-N_r-(1+m/2+2m^2)N_r=\frac{m^3}{2}H_7(m),
\]
where
\[
 H_7(m)=4+m(164-249m)+m^3(312-238m)
            +m^5(136-45m)+12m^7>0.
\]
Finally $1+m/2+2m^2-\ell(m)=2(m-1/2)^2$.

\Fcertificateheading{One residual and its exact polynomial definition}
Let $X=A$, $Y=\mathsf R^{-1}C$, and set
\[
 N_X=\|X\|^2,\quad N_Y=\|Y\|^2,\quad
 \nu=\langle X,Y\rangle,\quad
 \Delta_0=\frac{\operatorname{cross}(Y,X)}h.
\]
The ordered-chain geometry makes $\Delta_0>0$.  The Gram identity is
\begin{equation}
 h^2\Delta_0^2=N_XN_Y-\nu^2.
 \label{eq:certificate-Gram}
\end{equation}
For $t=1-2e$ define
\begin{equation}
 Q_X(e)=\Delta_0^2-\|tX-eY\|^2,\qquad
 Q_Y(e)=\Delta_0^2-\|tY-eX\|^2.
 \label{eq:certificate-Q}
\end{equation}
These are tangent residuals, not the frontier points $Q_j$.

Use barred coefficients $\bar\alpha=\alpha/h$, and similarly for the
other three frontier coefficients.  Put
\[
 A_1=\rho\bar\alpha,\quad B_1=\rho\bar\beta,\quad
 D_1=\rho\bar\delta,\quad G_1=\rho\bar\gamma,
 \qquad \rho=a^2+ab+b^2,
\]
\[
 n_b=12b^2-12b+7,\quad n_a=12a^2-12a+7,
 \qquad E_b=8\rho-12UA_1,\quad E_a=8\rho-12UD_1,
\]
\[
 \begin{aligned}
 J_b&=(1-2b)E_b+(A_1+2B_1)n_b,& L_b&=E_b-A_1n_b,\\
 J_a&=(1-2a)E_a+(D_1+2G_1)n_a,& L_a&=E_a-D_1n_a.
 \end{aligned}
\]
The denominators satisfy $E_a,E_b>6\rho>0$.
In an orthonormal frame the two vectors are
\[
 X=(-J_b/(2E_b),-hL_b/E_b),\qquad
 Y=(-J_a/(2E_a),hL_a/E_a).
\]
Consequently $4E_a^2E_b^2Q_X(e)$ equals
\[
 \begin{aligned}
 \mathcal N(D,e)={}&(J_bL_a+J_aL_b)^2
 -((1-2e)J_bE_a-eJ_aE_b)^2\\
 &-3((1-2e)L_bE_a+eL_aE_b)^2.
 \end{aligned}
\]
Reducing modulo $D^2=K:=4\rho-3$ defines, by exact polynomial division,
\begin{equation}
 \operatorname{rem}\mathcal N=2\rho^2(\mathcal A+D\mathcal B),\qquad
 \Pi=(1+3K)\mathcal A+K(3+K)\mathcal B.
 \label{eq:f-explicit-polynomials}
\end{equation}
Both $\mathcal B$ and $\Pi$ are quadratic in $e$.  Thus
\[
 Q_X(e)=\frac{\rho^2}{2E_a^2E_b^2}(\mathcal A+D\mathcal B).
\]
No independently stored coefficient table defines these polynomials.

\begin{lemma}[Radius transfer to the actual tangencies]
\label{lem:paper-residual-to-tangency}
If $0<e_0\le e_*<1/3$ and $Q_X(e_0),Q_Y(e_0)\ge0$, then the two
actual-radius tangent contact sums in
\zcref{eq:four-contact-tangent-tests} are at least $h$.
\end{lemma}
\begin{proof}
Apply \zcref{lem:paired-radius-transfer} to $X,Y,\Delta_0$.
This gives $Q_X(e_*),Q_Y(e_*)\ge0$.  The Gram identity gives
\begin{equation}
 P_A(he)=h^2N_XQ_X(e).
 \label{eq:certificate-PA}
\end{equation}
\begin{equation}
 P_C(he)=h^2N_YQ_Y(e).
 \label{eq:certificate-PC}
\end{equation}
The positive determinant sign and the unsquared implication after
\zcref{eq:four-contact-residuals} give the two support bounds.
Only the residual is tested at $e_0$; the disk in the line-contact and
exposed-hull arguments remains the actual disk of radius $he_*$.
\end{proof}

\Fcertificateheading{Explicit low-degree comparisons}
For $s,t\in[-1,1]$ write $m=(1+s)/4$, $U=(1+t)/12$.
The elementary Chebyshev recurrence
\[
 T_0(z)=1,\quad T_1(z)=z,\quad T_{n+1}(z)=2zT_n(z)-T_{n-1}(z)
\]
gives $|T_n(z)|\le1$ on $[-1,1]$.
If $F-g=\sum c_{ij}T_i(s)T_j(t)$, then
\begin{equation}
 |F-g|\le\sum_{i,j}|c_{ij}|.
 \label{eq:f-chebyshev-budget}
\end{equation}
The coefficients are computed over $\mathbb Q$, using the exact power
conversion and then checking reconstruction.  Thus each displayed error
budget below is a finite rational inequality, not a sampling claim.

\begin{lemma}[Explicit algebraic comparison bounds]
\label{lem:f-explicit-comparison-bounds}
On $\mathcal R=[0,1/2]\times[0,1/6]$ the polynomials defined by
\zcref{eq:f-explicit-polynomials} satisfy
\[
 \mathcal B(m,U,b_0)>4,\quad \mathcal B(m,U,m)>210,
 \quad \mathcal B_{ee}<-5600.
\]
Put
\[
 \begin{aligned}
 H_1&=-\Pi_{Uee},&
 H_2&=\Pi_{UUe}(m,U,0)-2\Pi_{Uee},\\
 H_3&=\Pi_{UUe}(m,U,1/3)-2\Pi_{Uee},&
 H_4&=\Pi_{Ue}(m,U,1/3).
 \end{aligned}
\]
Then $H_1>250000$, $H_2>500000$, $H_3>550000$, and $H_4>20000$.
Moreover
\[
 -\Pi_{UU}(m,U,1/3)+2\ell\Pi_{Ue}(m,U,1/3)-\ell^2\Pi_{ee}>21500.
\]
Writing $x=1-2m$, one also has
\[
 \begin{split}
 \Pi_U(m,U,b_0)\ge{}&140+503x+\frac{6265}{2}x^2\\
 &+25000(U-b_0)\qquad(U\ge b_0),
 \end{split}
\]
\[
 \Pi_U(m,U,1/3)>10700\qquad(m\ge1/3).
\]
On the strict physical domain, $N_X\ge N_Y$ and $\Pi_{ee}<0$.
The latter sign also holds on $0<U\le r(m)$.
\end{lemma}
\begin{proof}
We specify the certificate rather than an existence assertion for a sum
of squares.  The four normalized targets $H_i/1000$ have the quadratic
models
\[
 \begin{aligned}
 G_1={}&130s^2+80st-124s+52t^2+100t+454,\\
 G_2={}&466s^2+414st+124s+238t^2+771t+1379,\\
 G_3={}&286s^2+228st-202s+154t^2+348t+1090,\\
 G_4={}&18s^2+4st-14s+6t^2+17t+44.
 \end{aligned}
\]
Their exact Chebyshev error sums are bounded, respectively, by
$70,300,176,8$.  For $G=as^2+bst+cs+dt^2+jt+k$ complete the square as
\[
 G=a\left(s+\frac{bt+c}{2a}\right)^2
       +\gamma(t+1)^2+\delta(t+1)+L,
\]
where
\[
 \gamma=d-\frac{b^2}{4a},\quad
 \delta=j-\frac{bc}{2a}-2\gamma,\quad
 L=k-j+d-\frac{(c-b)^2}{4a}.
\]
For the four models $\gamma,\delta>0$, and
\[
 L=\frac{21188}{65},\quad\frac{373211}{466},\quad
   \frac{210031}{286},\quad\frac{57}{2},
\]
respectively.  Subtracting the error budgets proves the four bounds.

For $\mathcal B(m,U,b_0)$ use the cubic
\[
 \begin{split}
 10g={}&-292s^3+556s^2t+1504s^2+364st^2-317st-1795s\\
       &-28t^3+272t^2-389t+804.
 \end{split}
\]
Its exact error is less than $17/2$.  On $s\le0$, $g_{ss}>0$ and
$g_s(0,t)<0$, so $g(s,t)\ge g(0,t)\ge387/10$.
On $0\le s\le1$, the affine matrix $10(\nabla^2g-8I)$ has positive
upper-left entry and determinant at all four corners.  Its determinant
values are $55687,1026919,82551,23383$.  Convexity of the
positive-semidefinite cone gives $\nabla^2g\succeq8I$.
At $(2/3,1/3)$ the value is $383/30$ and gradient is
$(13/45,-35/18)$.  Completing a square gives
$g\ge1623259/129600>25/2$, hence $\mathcal B(m,U,b_0)>4$.
The other two $\mathcal B$ bounds follow by exact coefficient grouping
in $x=1-2m\in[0,1]$, $6U\in[0,1]$.

For the curvature target divided by $100$, use
\[
 \begin{split}
 g_c={}&76s^3-160s^2t-184s^2+104st^2+226st-132s\\
       &-8t^3-34t^2+283t+788.
 \end{split}
\]
The exact Chebyshev error is less than $55$.
Setting $X_0=1-s$, $Y_0=1+t$ and grouping nonnegative products on
$[0,2]^2$ gives
\[
 g_c\ge277-37X_0+52X_0^2
 =270+52(X_0-37/104)^2+87/208>270.
\]
This proves the curvature bound.  The first-derivative comparison uses
the polynomial increment
\[
 R(m,U)=\frac{\Pi_U(m,U,b_0)-\Pi_U(m,0,b_0)}{U},
\]
with its polynomial value at zero.  Two fixed-chart quadratic lower
models prove $R>25000$; the exact intercept grouping gives
$\Pi_U(m,0,b_0)\ge-2985+503x+12515x^2/2$.
Substituting $b_0=(1-x^2)/8$ proves the stated restricted estimate.
Grouping the $U$ coefficients for $m\ge1/3$ proves the bound $10700$.

For clarity, the two physical-domain signs are not claimed on an
unnecessary larger algebraic domain.  Exact norm subtraction factors as
\[
 N_X-N_Y=\frac{8U(a-b)\rho^2}{E_a^2E_b^2}\mathcal H(D,U).
\]
The explicit power-grouping proof gives
$\mathcal H\ge1261/162$ on $0\le U\le D^2/6$, $0\le D\le1$.
Also the definition of $\Pi$ has the weighted-conjugate identity
\[
 \Pi=\frac{(1+D)^3\mathcal N(D,e)+(1-D)^3\mathcal N(-D,e)}{4\rho^2}.
\]
Each numerator is a constant minus two squares of affine functions of
$e$, so $\Pi_{ee}<0$ for the physical $0<D<1$.
On $0<U\le r$, one has $K(m,U)<1$: at $c=C_L(m)$,
$K(m,1-c)-1=4(c-1)(c^3+c^2+c-2)<0$, and $K$ increases in $U$.

All expansions, power groupings, model reconstruction identities and
rational error bounds just specified are reproduced by the accompanying
\nolinkurl{verify_explicit_stage.py}.  The complete identities and
coefficient bounds, including the two increment models, are printed in
its colocated mathematical supplement.  This verifier uses no optimizer,
no stored SOS witness, no subdivision search, and no floating-point sign
test.  These finite coefficient computations are part of the proof;
they do not audit the preceding geometric forcing lemmas.
\end{proof}

\begin{lemma}[Boundary estimates and coordinatewise monotonicity]
\label{lem:f-boundary-comparisons}
For $0\le m\le1/2$,
\[
 \Pi(m,r,r)>\frac{1063}{1700}.
\]
The affine-branch left endpoints obey
\[
 \Pi(m,0,m)>0\quad(m\le1/3),
\]
\[
 \Pi(m,U,1/3)\ge4096/27
 \quad\bigl(m\ge1/3,\ U\ge(m-1/3)/3\bigr).
\]
Moreover $\Pi_e\ge1152$ on $\mathcal R\times[0,1/3]$.
\end{lemma}
\begin{proof}
At $U=0$, $x=1-2m$, exact factorization gives
\[
 \Pi_e(m,0,1/3)=\frac83(x^2+3)
 (21x^8+130x^6+289x^4+712x^2+144)\ge1152,
\]
and $\Pi_{ee}(m,0,e)<0$.  The positive $\Pi_{Ue}$ bound in
\zcref{lem:f-explicit-comparison-bounds} extends this to the whole
rectangle.  In particular $\Pi$ is increasing in both $U,e$ when
$U\ge b_0$ and $b_0\le e\le1/3$.

Put $q(m)=-9m^2/10+98m/125-3711/100000$.
The exact degree-eight comparison proves $b_0\le q\le r$ for
$1/5\le m\le1/2$.  Thus
\[
 \Pi(m,r,r)\ge
 \begin{cases}
 \Pi(m,b_0,b_0)>103,&m\le1/5,\\
 \Pi(m,q,q)>1063/1700,&m\ge1/5.
 \end{cases}
\]
The second polynomial has degree $27$.  With $x=m-41/100$, its explicit
Taylor coefficient bounds yield
\[
 \Pi(m,q,q)>16/25-10|x|+1700x^2\ge1063/1700.
\]
The first interval and the comparison $q\le r$ also use exact rational
coefficient bounds, reproduced in the same standalone verifier and
printed supplement.  They are not consequences of numerical fitting.
The parameter $q$ is only an auxiliary boundary comparison; it does not
replace $r$, $e_0$, or the actual disk.

For $m\le1/3$, the factor $(2m-1)^2$ can be removed from
$\Pi(m,0,m)$, and grouping higher negative powers under lower positive
powers on $0\le1-3m\le1$ proves positivity.
For $m\ge1/3$, put $x=1-2m\in[0,1/3]$.  Direct expansion gives
$\Pi(m,0,1/3)\ge-96+6688x^2/3$.  Integrating
$\Pi_U(m,U,1/3)>10700$ to $U\ge(1-3x)/18$ gives a decreasing
quadratic whose minimum is $4096/27$.
\end{proof}

\begin{theorem}[Explicit disk-tangent certificate]
\label{thm:paper-exact-mixed-certificate}
On the strict F domain, if $e_0\le1/3$, then
\[
 Q_X(e_0)>0,\qquad Q_Y(e_0)>0.
\]
The assertion uses the fixed-start comparison points, not the old
junction-start Newton points.
\end{theorem}
\begin{proof}
Concavity of $\mathcal B$ and its endpoint bounds give
$\mathcal B(m,U,e)>4$ on $b_0\le e\le m$, hence at $e=e_0$.
The elementary identity
\[
 D(1+3D^2)-D^2(3+D^2)=D(1-D)^3\ge0
\]
therefore gives
\[
 \mathcal A+D\mathcal B\ge\frac{\Pi}{1+3K}.
\]
It suffices to prove $\Pi(m,U,e_0)>0$.

For $U\ge r$, $e_0=r$; the restricted derivative estimate and
$\Pi_{Ue}>0$ give $\Pi_U(m,U,r)>0$.  The boundary estimate at $U=r$
settles this range.
For $U\le r$, put $f(U)=\Pi(m,U,m-\kappa U)$ and
$S_k(e)=\Pi_{UU}-2k\Pi_{Ue}+k^2\Pi_{ee}$.
For $0\le e\le1/3$ and the actual $\kappa\ge1$,
\[
 \partial_eS_\kappa=(1-3e)H_2+3eH_3+2(\kappa-1)H_1>0.
\]
On the integration path, $\Pi_{ee}<0$ and $\Pi_{Ue}>0$, so
$\partial_kS_k<0$ for $k\ge0$.  Consequently
\[
 f''(U)=S_\kappa(e)\le S_\kappa(1/3)
 \le S_{\ell(m)}(1/3)<-21500.
\]
The right endpoint is $U=r$.  The left endpoint is zero if $m\le1/3$;
otherwise it is $U_3=(m-1/3)/\kappa\ge(m-1/3)/3$.
All endpoint values are positive by
\zcref{lem:f-boundary-comparisons}.  Concavity proves positivity
throughout the interval.

Thus $Q_X(e_0)>0$.  Exact subtraction gives
\[
 Q_Y(e)-Q_X(e)=(1-e)(1-3e)(N_X-N_Y)\ge0
 \qquad(0\le e\le1/3),
\]
using the established norm ordering.  This proves both signs.
\end{proof}
